\documentclass[11pt]{article}

\usepackage{amsmath,amssymb,amstext,amsthm,mathtools}
\usepackage{geometry}
\usepackage{fancyhdr}
\usepackage[pdftex]{graphicx}
\usepackage{xcolor}

\geometry{
  verbose,
  dvips,
  width=430pt, marginparsep=0pt, marginparwidth=0pt,
  top=90pt, headheight=12pt, headsep=20pt, footskip=30pt, bottom=80pt}

\setlength{\parskip}{\medskipamount}
\setlength{\parindent}{0pt}

\linespread{1.05}

\newtheorem{theorem}{Theorem}
\newtheorem{lemma}[theorem]{Lemma}
\newtheorem{proposition}[theorem]{Proposition}
\newtheorem{corollary}[theorem]{Corollary}
\newtheorem{conjecture}[theorem]{Conjecture}
\newtheorem{obs}{Observation}
\newtheorem{clm}{Claim}
\newtheorem{ques}{Question}
\newtheorem{problem}{Problem}

\newtheorem{ex}{Example}


\newcommand{\spc}{\hspace{0.08in}}
\newcommand{\vs}{\vspace*{.1in}}
\newcommand{\E}{\mathbb{E}}
\newcommand{\Prob}{\mathbb{P}}
\newcommand{\edit}[1]{\emph{\textcolor{blue}{#1}}}
\newcommand*\dif{\mathop{}\!\mathrm{d}}



\renewenvironment{proof}{{\noindent \textbf{\textit{Proof.}}}}
{\hfill $\Box$ \vspace*{0.1in}}
\newenvironment{proofc}{{\noindent \textit{Proof of Claim.}}}
{\hfill $\Box$}





\begin{document}

Name: \underline{\hspace{10cm}}

\begin{center}
\begin{huge}\bf Multiple Integrals Practice 
\end{huge}
\end{center}

(1.) Consider the following expressions.

\begin{equation}
    I = \iint_R f(x,y) \dif{A} \hspace{1.5cm} \text{and} \hspace{1.5cm} J = \iiint_T f(x,y,z) \dif{V} 
\end{equation}



(i.) In words, what does $I$ represent?



\vspace{3cm}


(i.) In words, what does $J$ represent? (The 4D meaning)



\vspace{3cm}

(ii.) As a mathematical expression, what is the definition of $I$


\vspace{3cm}



(iv.) If $f(x,y,z)$ is a mass density function on the region $T$, what does $J$ represent?


\newpage

(2.) Consider the following double integral

\begin{equation}
    \iint_R 5x^3 \cos{y^3} \dif{A}
\end{equation}



where $R$ is bounded by $y=2$ and $y= \frac{1}{4}x^2$.

\vspace{.5cm}

(i.) Evaluate the double integral. 

\vspace{12cm}

(ii.) In words, interpret your previous answer.


\newpage

(3.) Consider the following double integral

\begin{equation}
    \iint_R e^{-x^2-y^2} \dif{A}
\end{equation}



where $R$ is the region between $x^2+y^2 = 2x$ and $x^2+y^2 = 8x$ and $y=0$.

\vspace{.5cm}

(i.) Evaluate the double integral. 

\vspace{11cm}

(ii.) Using calculus find the area of $R$.


\newpage


(4.) Consider the following triple integral.

\begin{equation}
    \iiint_T \sqrt{x^2+y^2} \dif{V}
\end{equation}

Where $T$ is the region which lies inside the cylinder $x^2+y^2 =16$, and between planes $z=1$ and $z = y+1$


(i.) Evaluate the triple integral.


\newpage

(5) Suppose $\rho(x,y,z) = 15z$ is the mass density function over the 3-dimensional region $T$ which is bounded by $2x+y+z=4$, $4x+4y+2z=20$, $z=2y^2$, and $z=\sqrt{4y}$.

(i.) Find the mass of the $T$.

\end{document}
